%=============================================================================!
% 10.7  MAKING IT FAST                                                        !
%=============================================================================!
\section{Making it fast}
\label{sec:md-speed}

A run of a million steps calls the model a million times, and the time a
call takes decides what can be simulated. This section measures where the
time goes and rewrites one force loop in three faster ways, each checked
against the plain one.

\subsection*{The profile}

A profiler records how long a program spends in each of its functions.
Python's \texttt{cProfile}, run on 200 steps of 4000 Lennard-Jones atoms
in reduced units with the NumPy pair model of \cref{sec:md-loop}, finds
three quarters of the time inside the neighbour list: 36\% in its 14
rebuilds by the cell list, whose loops over small cells are written in
Python, and 38\% in picking out at every step the listed pairs now within
$r_{\mathrm c}$, which reduces each of them to its nearest copy by solving
$\vb{h}\vb{s} = \vb{d}$ although the cell is a cube. The forces take most
of the remaining quarter, \texttt{np.add.at} alone 16\%. The first gains
would come from the neighbour list.

\subsection*{One loop, four ways}

The force loop runs over every pair at every step, and it shows the
choices well. A processor runs only machine code, its own short list of
simple instructions. A compiler can translate a numerical loop into machine code before it
runs. In the ordinary Python loop, the interpreter instead carries out
Python operations at each iteration, including the work needed to handle
objects and call NumPy functions. The kernel takes the
positions, the side lengths of a rectangular box and a list of pairs, and
returns the Lennard-Jones energy, shifted to zero at $r_{\mathrm c}$, and
the forces. \texttt{mdlab.kernels} has it four times:
\begin{itemize}
\item \texttt{lj\_loop}, a plain Python loop over the pairs;
\item \texttt{lj\_numpy}, the same arithmetic on whole arrays, every pair
   in each operation;
\item \texttt{lj\_numba}, the same loop with the three components of
   each pair written as a loop of their own, compiled to machine code by
   Numba when it is first called;
\item \texttt{lj\_fortran}, the loop written in Fortran and compiled
   through NumPy's \texttt{f2py}, which builds the code that lets Python
   call it.
\end{itemize}
The Fortran loop over the pairs is

\begin{code}[language={[95]Fortran}]
do p = 1, n_pairs
   i   = pair_i(p) + 1
   j   = pair_j(p) + 1
   sep = positions(:, j) - positions(:, i)
   sep = sep - box * anint(sep / box)
   r2  = sum(sep * sep)
   if (r2 >= r_cut2) cycle
   s6           = (sigma2 / r2)**3
   push         = 24.0_dp * well_depth * (2.0_dp * s6 * s6 - s6) / r2
   energy       = energy + 4.0_dp * well_depth * (s6 * s6 - s6)        &
                  - phi_cut
   forces(:, i) = forces(:, i) - push * sep
   forces(:, j) = forces(:, j) + push * sep
end do
\end{code}

\noindent Fortran counts from 1 where Python counts from 0, hence the
$+1$ of lines 2 and 3; \texttt{anint} rounds to the nearest whole number,
the minimum image of \cref{sec:md-image} in a rectangular box; and the
array \texttt{positions(3, n\_atoms)} holds one atom per column. Fortran
stores an array column after column and NumPy row after row, so both keep
the numbers in memory in the order $x_1, y_1, z_1, x_2, \dots$, and NumPy's
array of one atom per row, passed as its transpose, reaches Fortran
without a copy. \texttt{push} is $-\varphi'(r)/r$, so the force on $j$ is
\texttt{push} times the separation and the force on $i$ its opposite;
\texttt{push} times the separation is thus the opposite of the
\texttt{push} of \texttt{PairModel} in \cref{sec:md-loop}, which was the
force on $i$. The Numba version follows the Python loop with each
operation on three numbers written as a loop over the components, as in
the Fortran.

\Cref{fig:md-speed} times one call of each on atoms on a simple cubic
grid shaken by $0.05\sigma$, at $\rho = 0.8\sigma^{-3}$, with the pairs
within $r_{\mathrm c} + 0.3\sigma$ listed.
At $N = 2000$, with 74\,684 listed pairs, the plain loop takes
\SI{172}{\milli\second}, NumPy \SI{5.2}{\milli\second}, Numba
\SI{0.34}{\milli\second} and Fortran \SI{0.28}{\milli\second}: NumPy is 33
times faster than the loop, and the two compiled versions another 15 and
18 times faster than NumPy. Every version grows roughly as $N$, with the
number of listed pairs, the compiled ones less evenly. All three agree with the plain loop to about a part in $10^{12}$ in
the energy and to $\num{3e-12}$ in forces as large as 1024, in reduced
units, the differences of rounding in a different order.

\begin{figure}[htb]
\centering
\includegraphics{ch10_code/speed}
\caption{The time of one call of the Lennard-Jones force kernel against
$N$, for atoms on a simple cubic grid shaken by $0.05\sigma$ at $\rho =
0.8\sigma^{-3}$, with $r_{\mathrm c} = 2.5\sigma$, written as a plain
Python loop, with NumPy, with Numba and in Fortran, with a line of slope 1
(dashed). Each time is the best of five, on one thread, one sequence of
instructions on one core of the processor.}
\label{fig:md-speed}
\end{figure}

The reasons are in how each runs. The plain loop takes
$\SI{172}{\milli\second}/74\,684 = \SI{2.3}{\micro\second}$ for each
pair, and the Fortran loop \SI{3.7}{\nano\second} for the same
arithmetic, so almost all of the plain loop's time goes on something other
than arithmetic: each pair makes about a dozen calls of NumPy on arrays of
three numbers, and every call first finds out what its operands are and
makes a new array for its result. NumPy on whole arrays makes those
decisions once per operation for all the pairs and does the arithmetic in
compiled code, but it passes over the pairs once for every operation,
writing a new array each time. A compiled loop keeps each pair's few
numbers in the processor's registers, its fastest storage, and touches
memory only for the indices, positions and forces.

\begin{practice}
The profile comes before any rewriting: here the slow part was the
neighbour list. Only the loop that dominates is rewritten, the plain
version is kept as the reference, and every fast version is checked
against it, as the tests of \texttt{mdlab} do. Numba compiles a Python
loop with few changes once every operation on a small array is written as
a loop over its numbers, as in \texttt{lj\_numba}; the Fortran loop took
18\% less time here.
\end{practice}

\begin{keyidea}
A profiler shows where the time goes, and in this code it was the
neighbour list. The force loop written in Python, with NumPy on each pair,
is hundreds of times slower than the loop compiled by Numba or written in
Fortran; NumPy on whole arrays lies between. Every fast version is checked
against the plain one to the rounding of the arithmetic.
\end{keyidea}

\notebook{10}{race the four kernels}

\begin{selfcheck}
At \SI{0.2}{\micro\second} per atom per call, how long does the force
kernel alone take for \SI{1}{\nano\second} of $10^4$ atoms with $\dt =
\SI{1}{\femto\second}$?
\end{selfcheck}
